\documentclass[10pt,a4paper]{amsart}
\usepackage[margin=19mm]{geometry}
\usepackage{amsmath,amssymb,mathtools}
\usepackage[scr=rsfs,cal=euler]{mathalfa}
\usepackage{xfrac}
\usepackage[unicode,hidelinks,bookmarks=false]{hyperref}
\newcommand{\ie}{i.e.\ }
\newcommand{\cf}{cf.\ }
\newcommand{\resp}{respectively\ }
\newcommand{\Subsection}{\subsection}
\providecommand{\nobreakdash}{\nobreak\hbox{-}}
\setlength{\emergencystretch}{2em}
\multlinegap0pt
\usepackage{bm}
\usepackage{bbm}
\usepackage{dsfont}
\usepackage{xspace}
\usepackage{cancel}
\usepackage{mathtools}
\usepackage{amsthm}
\makeatletter
\@ifundefined{theorem}{\newtheorem{theorem}{Theorem}}{}
\@ifundefined{lemma}{\newtheorem{lemma}[theorem]{Lemma}}{}
\makeatother
\DeclareMathOperator{\Ind}{Ind}
\DeclareMathOperator{\diag}{diag}
\newcommand{\CC}{\mathbb{C}} %complex numbers
\newcommand{\FF}{\mathbb{F}} %rational numbers
\newcommand{\RR}{\mathbb{R}} %real numbers
\newcommand{\W}{\mathcal{W}}
\newcommand{\rank}{\mathrm{rank}} %rank
\newcommand{\tp}{{\scriptscriptstyle\mathsf{T}}}
\title[The Cayley Transform on Representations]{The Cayley Transform on Representations}
\author{Jingyu Lu, Ke Ye}
\date{Source: arXiv:2411.02071v2 (CC BY 4.0)}
\begin{document}
\maketitle

\noindent\emph{Source and licence.} The Cayley Transform on Representations. Version arXiv:2411.02071v2 (CC BY 4.0), distributed under the Creative Commons Attribution 4.0 International licence, as recorded in the arXiv licence metadata for this exact version and shown on the abstract page https://arxiv.org/abs/2411.02071v2. Full rights notice: licenses/arxiv-2411.02071-CC-BY-4.0.txt.\par\smallskip

\noindent\textit{Editorial scope.} Counted unit: the Lemma whose source label is \texttt{prop: Gauss elimination} in the article (source0.tex lines 511--519, source label \texttt{prop: Gauss elimination}), delivered together with the article's own complete original proof of it (source0.tex lines 520--556, one \texttt{proof} environment of 37 lines). No mathematics is written, repaired or re-derived: statement and proof are the article's own text line for line. Required context retained uncounted: lem:rank-orbit (lines 491--493); lem:rank-orbit:eq1 (lines 496--499); lem:rank-orbit:eq2 (lines 503--506). Every definition, display and label of those lines is retained unchanged. Omitted from this package and not redistributed: 1--490, 494--495, 500--502, 507--510, 557--976. Nothing inside a retained excerpt is omitted or altered: every retained body is delivered exactly as the article's lines are written, so no omission receipt of any kind accompanies this package. Citations: the retained excerpts carry no citation command at all, so no bibliography entry of the article is transcribed and no reference is redistributed. Reference adaptation: none was needed and none was made; every reference inside the delivered unit resolves inside it, so no unresolved reference or citation is produced. Printed slips kept literal: no printed word, symbol, hypothesis or number of the counted statement or of its proof was altered. Layout and class: the article's own class is \texttt{amsart} with packages including geometry, amsfonts, amsmath, amsthm, amssymb, scrextend, comment, tikz, tikz-cd, pgfplots, mathtools, euscript, rsfso, multicol; the wrapper is the lane's pinned \texttt{amsart} editorial wrapper at 10pt on A4, which restates the theorem-like environments the retained text needs and adds the article's own macro definitions that the retained text needs (\texttt{\textbackslash CC}, \texttt{\textbackslash FF}, \texttt{\textbackslash Ind}, \texttt{\textbackslash RR}, \texttt{\textbackslash W}, \texttt{\textbackslash diag}, \texttt{\textbackslash rank}, \texttt{\textbackslash tp}), with extra wrapper packages (bm, bbm, dsfont, xspace, cancel, mathtools). The delivered numbering is the unit's own. Assets: no retained span contains a figure environment, a graphics-inclusion command, a PDF-page inclusion, a table or a TikZ picture; no image file is copied into this package, the package declares no asset dependency and no image file is redistributed with it. Licence: version arXiv:2411.02071v2 of this article is distributed under the Creative Commons Attribution 4.0 International licence, as recorded in the arXiv licence metadata for this exact version and shown on the abstract page https://arxiv.org/abs/2411.02071v2. A full rights notice is delivered at licenses/arxiv-2411.02071-CC-BY-4.0.txt. Characters: every retained excerpt is pure ASCII, so the delivered engine, XeLaTeX with its default Latin Modern text font, reports no missing character.
\begin{lemma}\label{lem:rank-orbit}
If $\pi: \mathfrak{g} \to \mathfrak{gl}(\mathbb{V})$ is an irreducible representation with the highest weight $\mu$ and $\dim_\FF (\pi(\mathfrak{h}))= \dim_\FF(\mathfrak{h})$,  then $\rank ( \mathcal{O}_\mu ) = \dim_\FF(\mathfrak{h})$ where $\mathcal{O}_\mu \coloneqq \{s(\mu): s\in \W \}$ and $\W$ is the Weyl group of $\mathfrak{g}$.
\end{lemma}

\begin{equation}\label{lem:rank-orbit:eq1}
A = [
\underbrace{\lambda_1,\dots, \lambda_1}_{\dim (\mathbb{V}_{\lambda_1}) \text{~copies}},\dots,  \underbrace{\lambda_m,\dots, \lambda_m}_{\dim (\mathbb{V}_{\lambda_m}) \text{~copies}} ] \in \Lambda^{\dim (\mathbb{V})}.
\end{equation}

\begin{equation}\label{lem:rank-orbit:eq2}
A(H) \coloneqq \diag \left(
\underbrace{\lambda_1(H),\dots, \lambda_1(H)}_{\dim (\mathbb{V}_{\lambda_1}) \text{~copies}},\dots,  \underbrace{\lambda_m(H),\dots, \lambda_m(H)}_{\dim (\mathbb{V}_{\lambda_m}) \text{~copies}} \right) \in \mathbb{C}^{\dim (\mathbb{V}) \times \dim (\mathbb{V})}.
\end{equation}

\begin{lemma}\label{prop: Gauss elimination}
If $\pi: \mathfrak{g} \to \mathfrak{gl}(\mathbb{V})$ is an irreducible representation such that $\widetilde{\pi} \left( S^3 \mathfrak{h}\right)\subseteq \pi(\mathfrak{h})$ and  $\dim_\FF( \mathfrak{h} ) = \dim_\FF(  \pi(\mathfrak{h}) ) \eqqcolon n$,  then there exist {linearly independent} $\omega_1,\dots,  \omega_n \in \Lambda$ such that 
\[
\{\omega_1,\dots,  \omega_n\} \subseteq 
\{\lambda\in \Lambda\setminus \{0\}: \mathbb{V}_{\lambda} \ne 0\}
\subseteq 
\{\pm \omega_1,\dots,  \pm \omega_n\}. 
\]
\end{lemma}

\begin{proof}
Suppose $\Delta^+ = \{\alpha_1,\dots,  \alpha_n\}$.  Let $H_j$ be the co-root of $\alpha_j \in \Delta^+$ for each $1 \le j \le n $.  Since $\{H_1,\dots,  H_n\}$ is a basis of $\mathfrak{h}$,  it has the dual basis $\{\omega'_1,\dots,  \omega'_n\}$.  Assume $\{\lambda\in \Lambda: \mathbb{V}_{\lambda} \ne 0\}= \{ \lambda_1,  \dots,  \lambda_m \}$.  By definition,  there exists column vectors $a_1,\dots,  a_m \in \mathbb{Z}^{n}$ such that 
\begin{equation}\label{prop: Gauss elimination:eq1}
A =  \omega'  A',  \quad \omega'\coloneqq  \begin{bmatrix}
\omega'_1 &
\cdots &
\omega'_n
\end{bmatrix},\quad A' \coloneqq [
\underbrace{a_1,  \cdots,  a_1}_{\dim(\mathbb{V}_1) \text{~copies}},  \cdots,  \underbrace{a_m,  \cdots,  a_m}_{\dim(\mathbb{V}_m) \text{~copies}} ]
\end{equation}
where $A$ is defined in \eqref{lem:rank-orbit:eq1}.  Clearly,  $\omega'(\mathfrak{h}) \subseteq \mathbb{C}^n$ consists of vectors $\omega'(H) \coloneqq \begin{bsmallmatrix} \omega'_1(H),\cdots,  \omega'_n(H)
\end{bsmallmatrix}$ for $H\in \mathfrak{h}$.  In fact,  we must have $\omega'(\mathfrak{h}) = \mathbb{C}^n$ as $\rank (\omega') = n$.  Let $A_0 \coloneqq [
a_1, \cdots,  a_m] \in \mathbb{Z}^{n \times m}$.
\begin{itemize}
\item[$\diamond$] $\FF=\CC$.  According to the proof of Lemma~\ref{lem:rank-orbit},  we have $\rank(A_0) = n$.  We may re-index $\lambda_j$'s such that the left $n\times n$ submatrix of $A_0$ has the full rank.  By Gaussian elimination,  we obtain a decomposition $A_0 = GB$ where $G \in \mathbb{Z}^{n\times n}$ and 
\[
B = \begin{bmatrix}
I_n & C 
\end{bmatrix} = \begin{bmatrix} 
1 & 0 & \cdots & 0 & c_{1,1} & \cdots & c_{1,m-n} \\
0 & 1 & \cdots & 0 & c_{2,1} & \cdots & c_{2,m-n} \\
\vdots & \vdots & \ddots & \vdots & \vdots & \ddots & \vdots \\
0 & 0 & \cdots & 1 & c_{n,1} & \cdots & c_{n,m-n}
\end{bmatrix} \in \mathbb{Q}^{n \times m}.
\]
It is clear that we can construct $B' \in \mathbb{Q}^{n \times \dim (\mathbb{V})}$ from $B$ such that $A' = G B'$.  We take $\omega \coloneqq \omega' G \in \Lambda^n$.  By \eqref{prop: Gauss elimination:eq1} we have $A(H) = \diag ( \omega'(H) A' ) =  \diag ( \omega(H) B' )$ for any $H\in \mathfrak{h}$.  Since $S^3 \left( \pi(\mathfrak{h}) \right) = \widetilde{\pi} \left( S^3 \mathfrak{h}\right)\subseteq \pi(\mathfrak{h})$,  \eqref{lem:rank-orbit:eq2} implies that $A(H)^2 A(H') \in \pi (\mathfrak{h})$ for any $H,H'\in \mathfrak{h}$.  Let $X_1,\dots,  X_n$ be the dual basis of $\omega_1,\dots  \omega_n$.  Then by evaluating $A(X_j)^2 A(X_k)$ we obtain 
\[
c_{js}^3 = c_{js},\quad c_{js}^2 c_{ks} = 0,\quad 1 \le s \le m-n,\; 1 \le j < k \le n.
\]
Therefore,  $c_{js} \in \{0,-1,1\}$ for all $1 \le s \le m-n$ and $1 \le j  \le n$.  Moreover,  for each $1 \le s \le m-n$ there is at most one $1 \le j \le n$ such that $c_{js} \ne 0$.  By \eqref{prop: Gauss elimination:eq1},  we derive that $A = \omega B'$ where $B' \in \{0,\pm 1\}^{n \times \dim (\mathbb{V})}$.  By construction,  we clearly have $\omega = (\lambda_1,\dots,  \lambda_n)^\tp$ and this completes the proof.
\item[$\diamond$] $\FF=\RR$.  We notice that $        \widetilde{\pi} \left( S^3 \mathfrak{h} \right) \subseteq  \pi \left( \mathfrak{h} \right) $ implies 
\[
\widetilde{\Ind(\pi)} \left( S^3 (\mathfrak{h}^{\CC}) \right) = \widetilde{\pi} \left( S^3 \mathfrak{h} \right)\otimes \mathbb{C} \subseteq \pi \left( \mathfrak{h} \right) \otimes \CC = \Ind( \pi)\left( \mathfrak{h}^{\CC}\right).
\]
Since weights of $\pi$ are the same as those of $\Ind(\pi)$,  we obtain the desired inclusions by the complex case.  \qedhere
\end{itemize}
\end{proof}

\end{document}
